\documentclass[12pt, a4paper]{article}

\usepackage{fontspec}
\usepackage{polyglossia}
\setmainlanguage{french}
\usepackage{amsmath, amssymb}
\usepackage{geometry}
\usepackage{booktabs}
\usepackage{listings}
\usepackage{xcolor}
\usepackage{hyperref}
\usepackage{parskip}
\usepackage{algorithm}
\usepackage{algpseudocode}
\usepackage{graphicx}

\geometry{margin=2.5cm}

\definecolor{codebg}{RGB}{245,245,245}
\definecolor{codegreen}{RGB}{0,128,0}
\definecolor{codegray}{RGB}{128,128,128}
\definecolor{codeblue}{RGB}{0,0,200}

\lstset{
  language=Python,
  backgroundcolor=\color{codebg},
  basicstyle=\ttfamily\small,
  keywordstyle=\color{codeblue}\bfseries,
  commentstyle=\color{codegreen},
  stringstyle=\color{orange},
  numberstyle=\tiny\color{codegray},
  numbers=left,
  stepnumber=1,
  frame=single,
  framerule=0.4pt,
  breaklines=true,
  showstringspaces=false,
  tabsize=4,
}

\title{
  \textbf{Concorde TSP Solver}\\
  \large Documentation algorithmique --- Travelling Salesman Problem
}
\author{Recherche Opérationnelle --- CESI}
\date{}

\begin{document}

\maketitle
\tableofcontents
\newpage

% =============================================================================
\section{Définition}

\textbf{Concorde} est le solveur TSP exact de référence absolue, développé par
Applegate, Bixby, Chvátal et Cook à partir des années 1990. Il résout le TSP
à l'optimalité prouvée via \textbf{branch-and-cut} avec la formulation DFJ,
enrichie de coupes spécialisées (comb, blossom) issues de la théorie des
couplages parfaits.

Concorde a résolu des instances jusqu'à \textbf{85\,900 villes} (instance
\texttt{pla85900} de TSPLIB, 2006) et produit des bornes inférieures de
qualité exceptionnelle sur le World TSP à 1,9 million de villes. C'est le
seul algorithme à avoir prouvé l'optimalité sur des instances de cette taille.

% =============================================================================
\section{Méthode}

\subsection{Variables de décision}

\begin{equation}
  x_{ij} \in \{0, 1\}, \quad \forall\, i \neq j
\end{equation}

$x_{ij} = 1$ si l'arc $(i,j)$ est emprunté dans le tour optimal, $0$ sinon.

\subsection{Fonction objectif}

\begin{equation}
  \min \sum_{i \neq j} c_{ij} \cdot x_{ij}
\end{equation}

\subsection{Contraintes de degré}

\begin{align}
  \sum_{j \neq i} x_{ij} &= 1 \quad \forall\, i \\
  \sum_{i \neq j} x_{ij} &= 1 \quad \forall\, j
\end{align}

\subsection{Contraintes DFJ (lazy)}

\begin{equation}
  \sum_{\substack{i \in S,\, j \in S \\ i \neq j}} x_{ij} \leq |S| - 1
  \qquad \forall\, S \subsetneq V,\; 2 \leq |S| \leq n-1
  \label{eq:dfj}
\end{equation}

Ajoutées à la volée : la coupe associée à $S$ n'est générée que lorsque
la solution courante la viole. Voir la documentation ILP pour la preuve
complète d'élimination des sous-tours.

\subsection{Coupes de Comb (Padberg \& Rinaldi, 1990)}

Une \textbf{comb} est une structure combinatoire composée d'un \emph{manche}
$H \subset V$ et de $t \geq 3$ \emph{dents} $T_1, \ldots, T_t$ (nombre impair),
satisfaisant :

\begin{itemize}
  \item $T_k \cap H \neq \emptyset$ et $T_k \setminus H \neq \emptyset$ pour tout $k$
  \item $T_k \cap T_\ell = \emptyset$ pour tout $k \neq \ell$
\end{itemize}

L'inégalité de comb correspondante est :

\begin{equation}
  \sum_{(i,j) \in \delta(H)} x_{ij}
  + \sum_{k=1}^{t} \sum_{(i,j) \in \delta(T_k)} x_{ij}
  \geq t + 1
  \label{eq:comb}
\end{equation}

où $\delta(S)$ désigne les arcs traversant la frontière de $S$.

\textbf{Intuition :} tout tour hamiltonien doit traverser la frontière de $H$
un nombre pair de fois, et chaque dent au moins deux fois. Avec $t$ dents
impaires, le tour doit faire au minimum $t+1$ traversées au total.

\subsection{Coupes de Blossom (Edmonds, 1965)}

Les \textbf{blossom inequalities} sont des généralisations des inégalités de
comb issues de la théorie des couplages parfaits dans les graphes. Pour un
ensemble de nœuds $S$ avec $|S|$ impair :

\begin{equation}
  \sum_{(i,j) \in \delta(S)} x_{ij} \geq 2
  \label{eq:blossom}
\end{equation}

Plus généralement, les coupes blossom d'ordre supérieur impliquent des
structures en fleur (blossom) avec tige et pétales. Elles sont extrêmement
puissantes mais leur séparation (trouver laquelle est violée) est un problème
en soi, résolu par Concorde via des algorithmes de couplage minimum.

\subsection{Relaxation LP et qualité des bornes}

La force de Concorde réside dans la \textbf{qualité exceptionnelle de sa
relaxation LP}. En ajoutant des coupes DFJ, comb et blossom, la borne LP
converge très rapidement vers l'optimum entier. Sur la plupart des instances
pratiques, l'écart (gap) entre la borne LP finale et le tour optimal est
\textbf{inférieur à 1\%}, parfois nul --- ce qui signifie que l'optimalité
est prouvée sans aucun branchement.

% =============================================================================
\section{Pipeline de Concorde}

\begin{algorithm}[H]
\caption{Pipeline Concorde (branch-and-cut TSP)}
\begin{algorithmic}[1]
\State Initialiser avec les contraintes de degré uniquement
\State $z^* \leftarrow +\infty$ \Comment{meilleure solution entière connue}
\State File de nœuds $\mathcal{N} \leftarrow \{\text{racine}\}$
\While{$\mathcal{N} \neq \emptyset$}
  \State Extraire un nœud $P$ de $\mathcal{N}$
  \Repeat \Comment{boucle de coupes}
    \State Résoudre la relaxation LP de $P$ → $x^*_P$, borne $z_P$
    \If{$z_P \geq z^*$} \textbf{élaguer} $P$ \EndIf
    \State Chercher des sous-tours violés → ajouter coupes DFJ
    \State Chercher des combs violées → ajouter coupes \eqref{eq:comb}
    \State Chercher des blossoms violés → ajouter coupes \eqref{eq:blossom}
  \Until{aucune coupe violée ou solution entière}
  \If{$x^*_P$ est entière et forme un tour unique}
    \If{$z_P < z^*$} $z^* \leftarrow z_P$ ; enregistrer le tour \EndIf
  \Else \Comment{branchement}
    \State Choisir une variable fractionnaire $x_{ij}$
    \State Ajouter $P \cup \{x_{ij}=0\}$ et $P \cup \{x_{ij}=1\}$ à $\mathcal{N}$
  \EndIf
\EndWhile
\State \Return tour optimal de coût $z^*$
\end{algorithmic}
\end{algorithm}

% =============================================================================
\section{Explication intuitive}

Concorde résout le TSP comme un détective qui construit des preuves
d'impossibilité :

\begin{enumerate}
  \item \textbf{Phase de coupes} : il démontre mathématiquement que certaines
    configurations ne peuvent pas appartenir à un tour optimal. Chaque coupe
    (DFJ, comb, blossom) est une telle preuve. Plus les preuves sont précises,
    moins il reste de cas à explorer.

  \item \textbf{Phase de branchement} : quand les preuves ne suffisent plus,
    il divise le problème en deux cas ($x_{ij}=0$ ou $1$) et abandonne les
    branches qui ne peuvent pas améliorer la meilleure solution connue.
\end{enumerate}

En pratique, les coupes sont si puissantes que Concorde prouve souvent
l'optimalité \textbf{sans jamais brancher} : la relaxation LP donne
directement la solution entière optimale, avec une borne inférieure qui
coïncide exactement avec le tour trouvé.

\medskip
\textbf{Comparaison avec ILP MTZ :}

\begin{center}
\begin{tabular}{lll}
  \toprule
  & \textbf{ILP MTZ} & \textbf{Concorde (DFJ + coupes)} \\
  \midrule
  Contraintes initiales & $O(n^2)$ & Degré seulement \\
  Anti sous-tours       & Variables $u_i$ & Coupes DFJ lazy \\
  Coupes avancées       & Aucune & Comb, blossom \\
  Relaxation LP         & Faible & Excellente \\
  Branchements typiques & Beaucoup & Très peu, souvent zéro \\
  Limite pratique       & $n \leq 50$ & $n \leq 100\,000$ \\
  \bottomrule
\end{tabular}
\end{center}

% =============================================================================
\section{Cas d'applications connus}

\begin{table}[h]
  \centering
  \begin{tabular}{ll}
    \toprule
    \textbf{Référence} & \textbf{Contribution} \\
    \midrule
    Applegate, Bixby, Chvátal, Cook (1998)  & Première version publique de Concorde \\
    \textit{The Traveling Salesman Problem} (2006) & Livre de référence ; résolution de \texttt{pla85900} \\
    Padberg \& Rinaldi (1990--1991) & Coupes comb et blossom dans le pipeline B\&C \\
    Edmonds (1965) & Théorie des couplages parfaits --- fondement des blossoms \\
    World TSP (en cours) & 1\,904\,711 villes --- borne à $< 1\%$ de l'optimal \\
    \bottomrule
  \end{tabular}
  \caption{Références principales}
\end{table}

% =============================================================================
\section{Exemple de code}

L'implémentation ci-dessous reproduit le cœur du pipeline Concorde :
relaxation LP avec le moteur SCIP d'OR-Tools, coupes DFJ lazy, et
branch-and-cut intégré au solveur MIP.

\begin{lstlisting}[caption={Concorde-inspired : LP + DFJ lazy + branch-and-cut (OR-Tools SCIP)}]
from ortools.linear_solver import pywraplp
import math

def find_subtours(n, x_val):
    visited, subtours = [False] * n, []
    for start in range(n):
        if visited[start]:
            continue
        tour, cur = [], start
        while not visited[cur]:
            visited[cur] = True
            tour.append(cur)
            nxt = next((j for j in range(n) if j != cur and x_val[cur][j] > 0.5), None)
            if nxt is None:
                break
            cur = nxt
        subtours.append(tour)
    return subtours


def solve_tsp_concorde(cities, time_limit=120.0):
    """
    Concorde-inspired solver : LP relaxation + DFJ lazy constraints
    + branch-and-cut via OR-Tools SCIP.
    Retourne : (tour, distance, nb_iterations, nb_coupes)
    """
    n = len(cities)

    def dist(i, j):
        dx = cities[i][0] - cities[j][0]
        dy = cities[i][1] - cities[j][1]
        return math.hypot(dx, dy)

    forbidden = []
    iteration = 0

    while True:
        iteration += 1
        solver = pywraplp.Solver.CreateSolver("SCIP")
        solver.SetTimeLimit(int(time_limit * 1000))

        x = [[solver.IntVar(0, 1, f"x_{i}_{j}") for j in range(n)]
             for i in range(n)]

        # Contraintes de degre
        for i in range(n):
            solver.Add(x[i][i] == 0)
            solver.Add(sum(x[i][j] for j in range(n) if j != i) == 1)
            solver.Add(sum(x[j][i] for j in range(n) if j != i) == 1)

        # Coupes DFJ accumulees (lazy)
        for S in forbidden:
            S_list = list(S)
            solver.Add(
                sum(x[i][j] for i in S_list for j in S_list if i != j)
                <= len(S_list) - 1
            )

        # Objectif
        solver.Minimize(
            sum(dist(i, j) * x[i][j]
                for i in range(n) for j in range(n) if i != j)
        )

        status = solver.Solve()
        if status not in (pywraplp.Solver.OPTIMAL, pywraplp.Solver.FEASIBLE):
            return [], float("inf"), iteration, len(forbidden)

        x_val = [[x[i][j].solution_value() for j in range(n)]
                 for i in range(n)]
        subtours = find_subtours(n, x_val)

        if len(subtours) == 1:
            break

        for st in subtours:
            fs = frozenset(st)
            if fs not in forbidden and len(fs) < n:
                forbidden.append(fs)

    tour = subtours[0]
    total = (sum(dist(tour[k], tour[k+1]) for k in range(n-1))
             + dist(tour[-1], tour[0]))
    return tour, total, iteration, len(forbidden)
\end{lstlisting}

% =============================================================================
\section{Benchmark TSP}

Mesuré sur graphes aléatoires (seed = 42, coordonnées uniformes dans $[0,100]^2$),
solveur OR-Tools SCIP, limite de temps 60\,s par instance, CPU.

\begin{table}[h]
  \centering
  \resizebox{\textwidth}{!}{%
  \begin{tabular}{ccccccc}
    \toprule
    \textbf{Nœuds} & \textbf{\% réussite} & \textbf{\% non-détection}
    & \textbf{\% fausse détection} & \textbf{Temps CPU}
    & \textbf{Itérations} & \textbf{Statut} \\
    \midrule
    1        & ---   & --- & --- & $< 0.01$\,s & ---  & INFEASIBLE \\
    10       & /     & /   & /   & /           & /    & /          \\
    100      & /     & /   & /   & /           & /    & /          \\
    1\,000   & /     & /   & /   & /           & /    & /          \\
    10\,000  & /     & /   & /   & /           & /    & /          \\
    100\,000 & /     & /   & /   & /           & /    & /          \\
    \bottomrule
  \end{tabular}}
  \caption{Résultats benchmark Concorde-inspired (OR-Tools SCIP + DFJ lazy).}
\end{table}

\end{document}
